\documentclass[../main.tex]{subfiles}

\begin{document}

\chapter{Evaluasi Pernyataan Majemuk dengan Tabel Kebenaran}

\begin{subcpmk}
  \item Mahasiswa mampu menggunakan tabel kebenaran untuk mengevaluasi pernyataan majemuk yang kompleks.
  \item Mahasiswa dapat merunut dan mensimulasikan nilai aljabar logika untuk banyak variabel komputasi.
\end{subcpmk}

\section{Penggunaan Tabel Kebenaran (Truth Table)}
Setiap proposisi elementer dibatasi atas dua keadaan (Benar atau Salah). Ketika terdapat sekumpulan sejumlah $n$ variabel proposisi deklarasi ($p_1, p_2, \ldots p_n$), kombinasi evaluasi total probabilitas nilai kebenaran akan memenuhi perumusan ruang probabilitas $2^n$ baris data di dalam Tabel Kebenaran.
Tabel Kebenaran adalah peranti algoritmis grafis yang berdaya guna dalam mendeteksi dan mengeksekusi semantik Boolean secara terpusat dan tanpa ambiguitas operasional. Langkah iterasi ini serupa dengan algoritma pereduksi \textit{brute-force} guna me-list (\textit{enumeration}) probabilitas mutasi *state* / status dari \textit{software}.

\subsection{Tahapan Penyusunan Tabel}
Misalkan sebuah mesin komputer mengevaluasi urutan ekspresi komposit:
$X = (p \land \sim q) \lor r$
\begin{enumerate}
    \item Hitung iterasi $n$. Ada 3 buah var ($p, q, r$). Maka total baris wajib dicitrakan sebanyak $2^3 = 8$ baris konjugasi nilai (Kombinasi 000 hingga 111).
    \item Gambarkan kolom untuk setiap unit atom dasarnya terlebih dahulu (kolom atomik primer).
    \item Lanjutkan pembentukan kolom komponen secara sub-himpunan, yakni mengartikulasikan bentuk bersarang yang paling dalam tanda kurungnya dulu, contoh: $\sim q$, kemudian $(p \land \sim q)$.
    \item Bentuk kolom evaluasi kumulatif keseluruhan: $(p \land \sim q) \lor r$.
\end{enumerate}

\begin{contoh}
Tabel skematik untuk 2 var iterasi ($p, q$) membentuk 4 baris evaluasi:

\begin{center}
\begin{tabular}{|c|c|c|c|c|c|}
\hline
\textbf{$p$} & \textbf{$q$} & \textbf{$\sim q$} & \textbf{$p \land q$} & \textbf{$p \lor q$} & \textbf{$p \rightarrow \sim q$} \\
\hline
T & T & F & T & T & F \\
T & F & T & F & T & T \\
F & T & F & F & T & T \\
F & F & T & F & F & T \\
\hline
\end{tabular}\\
\textit{Contoh Tabel Kebenaran Sintesis $p \rightarrow \sim q$}
\end{center}
\end{contoh}

\begin{aktivitas}
  \item Simulasikanlah, bagaimana komputasi program jika ia menerima argumen Boolean yang lebih makroskopis, misalnya melibatkan susunan 5 variabel independen pengontrol navigasi pesawat (sebanyak $2^5 = 32$ permutasi). Apa kendala manusia yang diatasi melalui tabel kebenaran sistematis?
\end{aktivitas}

\begin{latihan}
  \item Susunlah Tabel Kebenaran (T-F / 1-0) iteratif bagi fungsi majemuk komputasi $f(A,B,C): A \leftrightarrow (B \land \sim C)$. Ada berapa permutasi \textit{assignment} baris yang menjadikannya *True/Benar* pada finalisasinya?
\end{latihan}

\begin{checklist}
  \item Memahami konsep probabilitas jumlah baris data $2^n$.
  \item Mampu mensintesis tabel dan membongkar sub-kolom matematis fungsi bersarang tanpa salah iterasi nilai.
\end{checklist}

\end{document}
